\begin{table}
\begin{tcolorbox}[tabularx={X|p{2cm}|X|},enhanced,fonttitle=\bfseries\large,fontupper=\normalsize\sffamily,
colback=yellow!10!white,colframe=red!50!black,colbacktitle=blue!30!white,
coltitle=black,title=Classifying Arcs]
\hline
Command  & Arguments & Purpose \\
\hline
\hline
\verb'-control' & label & Set the label of the poset classification control object. \\
\hline
\verb'-projective_space' & P & Set the projective space to P. \\
\hline
\verb'-d' & $d$ & Require that no more than $d$ points lie on a line. \\
\hline
\verb'-target_size' & $t$ & Specify the size of the arc to be $t$. \\
\hline
\verb'-conic_test' & & Require that no $6$ points of the arc lie on a conic. \\
\hline
\verb'-test_nb_Eckardt_points' & $n$ & Require exactly $n$ Eckardt points. \\
\hline
\verb'-affine' & & Classify arcs in the affine geometry, assuming that $x_0 = 0$ is the hyperplane at infinity. The condition that no more that $d$ point lie on a line applies to affine lines only. \\
\hline
\verb'-no_arc_testing' & & Do not test the ``at most $d$ points per line'' condition. \\
\hline
\verb'-forbidden_point_set' & set & The arc must not contain any of the given points. \\
\hline
\verb'-override_group' & label & Override the group under which we classify. \\
\hline
\end{tcolorbox}
\caption{\label{tab:arcs}Classifying Arcs}
\end{table}
